\section{Preentrenamiento en dos fases: primero expandir la cobertura, luego concentrar la calidad}

El preentrenamiento en dos fases no consiste simplemente en dividir el entrenamiento por la mitad, sino en utilizar dos distribuciones de muestreo con objetivos diferentes. La fase 1 expone al modelo a un conocimiento del mundo amplio y diversas expresiones; la fase 2, manteniendo las capacidades generales, aumenta la probabilidad de muestreo de fuentes matemáticas, de código, Wikipedia, de tareas y de otros orígenes de alta calidad. El orden mismo se convierte en una variable optimizable.

\subsection{El mecanismo central del currículo por tramos}

La palabra clave de la fase 1 es la diversidad: grandes cantidades de rastreo web (web crawl), menor repetición de datos de alta calidad y una cobertura de dominios más amplia. La palabra clave de la fase 2 es la calidad: aumentar el número de épocas con datos de alta calidad, utilizando las actualizaciones limitadas tardías para señales de capacidad más densas. Las áreas roja y verde en el gráfico no representan una división moral entre "datos malos" y "datos buenos", sino una asignación de recursos por etapas diferentes.

\lecturefigure{slide-010.jpg}{Preentrenamiento en dos fases: la fase 1 prioriza la diversidad y la fase 2 se centra en la calidad}{Grabación oficial de Stanford, 00:09:29--00:10:14}

El currículo puede escribirse como una distribución de muestreo por tramos:
\[
p(D_k\mid s)=
\begin{cases}
w_k^{(1)}, & 0\le s<s_0,\\
w_k^{(2)}, & s_0\le s\le 1,
\end{cases}
\qquad
w_{\text{HQ}}^{(2)}>w_{\text{HQ}}^{(1)}.
\]
Donde $s$ es el progreso de entrenamiento normalizado, $s_0$ es el límite entre fases (phase boundary), y $w_k^{(1)}$ y $w_k^{(2)}$ son los pesos de los datos para ambas etapas. El informe del artículo indica que utilizar la fase 2 en aproximadamente el último 40\% del entrenamiento produce mejores resultados en su configuración, pero $s_0$ aún debe determinarse mediante experimentos proxy.

\begin{lstlisting}[caption={Muestreador two-phase bajo presupuesto de token fijo}]
def sample_source(step, total_steps, phase_boundary, phase1, phase2):
    progress = step / total_steps
    weights = phase1 if progress < phase_boundary else phase2
    return categorical_sample(weights)

for step in range(total_steps):
    source = sample_source(step, total_steps, 0.60, p1_weights, p2_weights)
    batch = next_batch(source)
    optimize_next_token_loss(batch)
\end{lstlisting}

\subsection{Tres líneas base (baselines) que aíslan tres variables}

La distribución natural muestrea según el número de tokens existentes en cada fuente de datos: las fuentes grandes y de baja calidad ocuparán más actualizaciones debido a su escala. El *optimal blend* con orden aleatorio ya utiliza información de calidad y épocas, pero no altera el orden temporal. El enfoque de dos fases controla simultáneamente la calidad, las épocas y el orden. Solo con un diseño de líneas base así se puede separar el beneficio del "mejor mezcla" del del "mejor currículo".

\lecturefigure{slide-011.jpg}{Distribución natural, mezcla aleatoria óptima y línea base two-phase}{Grabación oficial de Stanford, 00:10:14--00:11:42}

\begin{knowledgebox}{Lectura de la figura: primero observe las tres columnas, luego los nombres del modelo}
Las tres columnas corresponden a Calidad (Quality), Épocas (Epochs) y Orden (Order). La distribución natural no controla ninguno de los tres; el *optimal random blend* controla los dos primeros; *two-phase* controla además el orden. Los nombres Pascal y Volta son solo ayudas mnemotécnicas. Si algún nuevo experimento modifica simultáneamente el tokenizador, la escala del modelo o el número total de tokens, la comparación ya no será limpia según la definición de esta tabla.
\end{knowledgebox}

\subsection{¿Qué responden el 17\% y el 3.4\% respectivamente?}

La sección anterior ya separó las variables de control entre las tres líneas base; esta sección explica con mayor detalle a qué contrafactual corresponde cada uno de los dos números destacados, evitando confundir el "beneficio total de optimizar la mezcla" con el "beneficio aportado únicamente por el orden". En el gráfico de resultados, la mejora relativa promedio de Volta frente a Pascal es de aproximadamente 17\%, donde tanto la mezcla de calidad como el orden cambian. Frente a la línea base que ya posee *optimal blend* pero con orden aleatorio, *two-phase* sigue mostrando una mejora de aproximadamente 3.4\%. El primer número indica la importancia de la receta (*recipe*), mientras que el segundo se acerca más al aporte independiente del orden del currículo.

\lecturefigure{slide-012.jpg}{Resultados de la estrategia two-phase frente a la distribución natural y la mezcla aleatoria óptima}{Grabación oficial de Stanford, 00:11:42--00:12:32}

Si la precisión de la línea base es $a$ y la del nuevo método es $b$, la mejora relativa es
\[
\Delta_{\mathrm{rel}}=\frac{b-a}{a}\times 100\%,
\]
mientras que la mejora absoluta es $\Delta_{\mathrm{abs}}=b-a$ puntos porcentuales. Las menciones de "\% mejor" en el curso se refieren a mejoras relativas y no deben mezclarse con puntos porcentuales absolutos, ni sumarse directamente entre diferentes suites de benchmarks.

\teachervoice{00:27:16--00:27:58, cuando el orador responde a si es factible "primero calidad, luego diversidad", dice que intercambiar las dos fases en un estudio de ablicación (ablation) da resultados peores. Este resultado apoya la dirección actual, pero sigue siendo una experiencia bajo un corpus y régimen de modelo específicos, no una ley universal del aprendizaje.}

\begin{warningbox}{Las dos fases no significan "primero basura, luego joyas"}
La fase 1 aún requiere filtrado y control básico de calidad; la fase 2 tampoco debe limitarse a retener solo datos de puntuación estrecha y alta. Acortar demasiado la cobertura causaría pérdida de alcance; repetir datos de alta calidad demasiado tarde o durante un periodo excesivo produciría rendimientos decrecientes (diminishing returns). En términos de ingeniería, se deben monitorear simultáneamente la cobertura de dominios, la pérdida de validación (*validation loss*), el fenómeno de memorización y la transferencia a tareas posteriores (downstream transfer).
\end{warningbox}

\subsection{Resumen de la sección}

El preentrenamiento en dos fases incorpora el currículo en la distribución de muestreo: primero se expande la cobertura y, posteriormente, se aumenta la densidad de señales de alta calidad. Una explicación rigurosa debe distinguir entre distribución natural, mezcla aleatoria óptima y mezcla ordenada, y especificar las métricas de mejora relativa. A continuación, el curso avanza desde los datos generales hacia los más sensibles: los datos de razonamiento (*reasoning data*).
